\chapter{Size Issues in Category Theory}

In an introduction to category theory, one often learns that a category $\catC$ consists of a \emph{collection} of objects (together with morphisms, composition operations, etc. satisfying some axioms). Often in mathematics, the word ``collection'' is synonymous with the word ``set'', but not here! What is meant precisely by the word ``collection'' in the definition of a category depends slightly on how one wishes to formalize their set-theoretic foundations for category theory, and there are generally two popular approaches here. We will work in ZFC set theory.

\section{Classes}

A famous argument of Russell shows that there is no set of all sets, i.e.

\begin{equation}\label{eq:russel}
    \nexists A \forall x (x \in A).
\end{equation}

We quickly recount the proof.

\begin{proof}
    Suppose (towards a contradiction) that the set $A$ of all sets did exist, i.e. $\exists A \forall x (x \in A)$. Then one could form the set
    \[B \coloneq \{x \in A : x \notin x\}.\]
    Now by construction (and since $B \in A$), we have $B \in B \iff B \notin B$, which is a contradiction.
\end{proof}

We emphasize that \eqref{eq:russel} is a statement in the language of set theory, and so its proof relies on some axioms of set theory. In this case, the proof of \eqref{eq:russel} relies on only one axiom (schema), the \emph{axiom schema of restricted comprehension}.

\begin{axiom}[Axiom Schema of Restricted Comprehension\footnote{This is actually weaker than the standard axiom schema of restricted comprehension used in (e.g.) ZFC set theory, but the details are unimportant here.}]
    Let $\varphi(x)$ be a formula in the language of set theory (i.e. $\varphi$ may have $x$ as a free variable and no other free variables).

    For all $A$, there exists $B$ such that for all $x$, $x \in B$ if and only if $x \in A$ and $\varphi(x)$ holds.
\end{axiom}

This is called an axiom \emph{schema} because it is actually an infinite collection of axioms (a different axiom for each formula $\varphi$). It may be tempting to try to rewrite this axiom schema as a single axiom of the form $\forall \varphi(\cdots)$, but this does not make sense -- axioms of set theory must be statements in the language of set theory. Formulas in the language of set theory are not themselves sets, so we may not quantify over formulas!

Despite its name, the axiom schema of restricted comprehension is easy to understand. In plain English: given any set $A$ and any proposition $\varphi(x)$, one may construct a set consisting precisely of those elements of $A$ which satisfy the proposition $\varphi$. If one also has the axiom of extensionality\footnote{The axiom of extensionality states that sets are determined by their elements, i.e. $\forall A \forall B (A = B \iff \forall x (x \in A \iff x \in B))$. We will assume this axiom always, though the proof of \eqref{eq:russel} given above does not actually require the axiom of extensionality.}, which one should always include in their axioms of set theory, then the set $B$ described above is necessarily unique, and we write
\[\{x \in A : \varphi(x)\}\]
to denote this set. Indeed, the axiom schema of restricted comprehension is what gives meaning to this set-builder notation! If you have ever defined a set like
\[\{x \in \mathbb{R} : \lvert x \rvert < 1\},\]
you were using the axiom schema of restricted comprehension. The tool of set-builder notation is crucial for developing almost any piece of mathematics, and so the axiom schema of restricted comprehension is essential in any useful theory of sets.

Now the key point is that the axiom schema of restricted comprehension allows us to build \emph{subsets} $\{x \in A : \varphi(x)\}$ of any \emph{pre-existing set} $A$, but it does \emph{not} allow us to construct new sets which are not subsets of a pre-existing set! In particular, the axiom schema does \textbf{not} allow us to interpret
\[\{x : x \notin x\}\]
(for example) as a well-defined set\footnote{This expression is an example of an \emph{unrestricted} comprehension, which our axiom says nothing about.}. Indeed, as we have seen, the axiom schema of restricted comprehension allows us to prove that no such set exists! On the other hand, for any \emph{given} set $A$, one \emph{may} form
\[A_\top \coloneq \{x \in A : x \notin x\},\]
i.e. the set of all elements of $A$ which do not contain themselves\footnote{Using the \emph{foundation} and \emph{pairing} axioms of ZFC, one may prove $\forall x (x \notin x)$, and conclude that $A_\top = A$ for any set $A$.}. This is simply the set whose existence is guaranteed by the axiom schema of restricted comprehension applied to the formula $\varphi(x)$ given by ``$x \notin x$''. For a more extreme example, one may not interpret
\[\{x : x = x\}\]
as a set (it would be the ``set of all sets'', which we have shown does not exist), but one may form the set $\{x \in A : x = x\}$ given any set $A$ (which will always be equal to $A$).

On the other hand, we would like to have access to the category of sets, which we will denote $\Set$. Somehow, every set should be an object of $\Set$, i.e. the collection of objects of $\Set$ should be the collection of all sets (whatever that means). We know that there is no set of all sets, so this ``collection of objects'' must not be a set, but instead something else.

\begin{definition}
    A \emph{class} is an equivalence class of formulas (in the language of set theory) with one free variable, with respect to logical equivalence. We say that a formula $\varphi(x)$ is logically equivalent to another formula $\psi(x)$ if and only if
    \[\forall x (\varphi(x) \iff \psi(x))\]
    holds. We may write $\{x : \varphi(x)\}$ to denote the equivalence class of a formula $\varphi(x)$.

    For any formula $\varphi(x)$ and any set $y$, we may write $y \in \{x : \varphi(x)\}$ to mean $\varphi(y)$, which we note is well-defined by the definition of logical equivalence of formulas.
\end{definition}

Now we have a way (tautologically) to interpret
\[\{x : x = x\}\]
as a \emph{class}, and we call this class the \emph{class of all sets}. Now it is easy to prove the statement
\[\forall y (y \in \{x : x  = x\})\]
since this is just shorthand for $\forall y (y = y)$, which is true! Simply put, classes are \emph{unrestricted} comprehensions, and they behave just like sets (for any given set $S$ and any given class $C$, the statement $S \in C$ is either true or false), except that they are not sets (e.g. one may not quantify over classes -- the variables appearing in a formula in the language of set theory are sets!). For example, we have a class $C = \{x : x \notin x\}$, but we cannot consider the statement ``$C \in C$'' to derive a contradiction: it does not make sense to ask if one class is an element of another! The elements of a class are sets.

Now consider the class $\{x : x \neq x\}$. One can just as easily prove that $\forall y (y \notin \{x : x \neq x\})$, since this is shorthand for $\forall y(\lnot (y \neq y))$, which is true. In other words, we should think of $\{x : x \neq x\}$ as the ``empty class''. There is also the empty set $\varnothing$\footnote{The existence of $\varnothing$ follows (in ZFC) from the axiom of infinity together with the axiom schema of restricted comprehension.}, and we note that
\[\forall y (y \in \varnothing \iff y \in \{x : x \neq x\})\]
holds. In other words, the class $\{x : x \neq x\}$ has the same elements as the set $\varnothing$! We should therefore formally identify the class $\{x : x \neq x\}$ with the set $\varnothing$, and in general identify any class $C$ with any set $S$ whenever $C$ and $S$ have the same elements.

\begin{convention}
    Let $C$ be a class and let $S$ be a set. If
    \[\forall x (x \in C \iff x \in S)\]
    holds, we identify $C$ with $S$ and (in particular) say that $C$ \emph{is a set}. If $C$ is not identified with any set $S$, i.e. if
    \[\forall S \exists x (\lnot(x \in C \iff x \in S))\]
    holds, then we say that $C$ is a \emph{proper class}.
\end{convention}

We note that this does not cause any additional identifications: if classes $C_1$ and $C_2$ are both identified with a set $S$, then $\forall x (x \in C_1 \iff x \in C_2)$ holds, so already $C_1 = C_2$. Likewise, if a class $C$ is identified both with a set $S_1$ and with a set $S_2$, then $\forall x (x \in S_1 \iff x \in S_2)$ holds, so already $S_1 = S_2$.

We have seen one example of a proper class already, namely the class of all sets $\{x : x = x\}$. There are many others, for example the class of all groups
\begin{equation}\label{eq:class of groups 1}
    \{x : \exists G \exists {\star} (x = (G,\star) \text{ and } (G,\star)\text{ is a group})\}.
\end{equation}
Of course, here the text ``$(G,\star)\text{ is a group}$'' stands in for a very long statement in the language of set theory which encodes the group axioms for the pair $(G, \star)$. We might also write (as shorthand)
\begin{equation}\label{eq:class of groups 2}
    \{(G,\star) : (G,\star)\text{ is a group}\}
\end{equation}
for this class, abusing notation since this is not precisely in the form $\{x : \varphi(x)\}$. It's not a big deal, because there is a canonical procedure to rewrite \eqref{eq:class of groups 2} in the standard form \eqref{eq:class of groups 1}.

In the other direction, we have seen an example of a non-proper class, i.e. a class which is also a set, namely $\{x : x \neq x\}$ (which is equal to $\varnothing$).

We may now make precise the definition of a category. Important for the definition will be the notion of a function between classes.

\begin{definition}
    Let $C$ and $C'$ be classes. We say $C \subseteq C'$ if and only if
    \[\forall x (x \in C \implies x \in C').\]
\end{definition}

Note that this extends the subset relation between sets, since every set is also (identified with) a class!

\begin{definition}
    Let $C$ and $C'$ be classes. We denote by $C \times C'$ the class
    \[\{(x,y) : x \in C \land y \in C'\},\]
    which (as a reminder) is shorthand for
    \[\{p : \exists x \exists y (x \in C \land y \in C' \land p = (x,y))\}.\]
\end{definition}

\begin{definition}
    Let $C$ and $C'$ be classes. By a \emph{function} $C \to C'$ we will mean a class $f \subseteq C \times C'$ such that
    \[\forall x (x \in C \implies \exists! y (y \in C' \land (x,y) \in f))\]
    holds.
\end{definition}

Note that when $C$ and $C'$ are sets, this recovers exactly the notion of a function between sets. In any case, when $f : C \to C'$ is a function between classes and $x \in C$ is an element, we denote by $f(x)$ the unique element of $C'$ such that $(x,f(x)) \in f$.

\begin{definition}[Categories, Version 1]
    A category $\catC$ consists of the following data:
    \begin{itemize}
        \item A class $\ob \catC$ of \emph{objects},
        \item A class $\mor \catC$ of \emph{morphisms},
        \item A function $\dom : \mor \catC \to \ob \catC$,
        \item A function $\cod : \mor \catC \to \ob \catC$,
        \item A function $\id : \ob \catC \to \mor \catC$,
        \item A function $\circ : \{(g,f) \in \mor \catC \times \mor \catC : \dom g = \cod f\} \to \mor \catC$;
    \end{itemize}
    satisfying the following axioms:
    \begin{enumerate}
        \item For all $x \in \ob \catC$, we have $\dom \id x = x = \cod \id x$,
        \item For all $f \in \mor \catC$, we have $f \circ \id \dom f = f = \id \cod f \circ f$,
        \item $\circ$ is associative, i.e. $f \circ (g \circ h) = (f \circ g) \circ h$ whenever this makes sense.
    \end{enumerate}
    In this description, we get for any two objects $x,y \in \ob \catC$ a \emph{hom-class}
    \[\Hom_{\catC}(x,y) \coloneq \{f \in \mor \catC : \dom f = x \land \cod f = y\}.\]
    If $\Hom_{\catC}(x,y)$ is a set for all $x, y \in \ob \catC$, we say that $\catC$ is \emph{locally small}. If $\mor \catC$ is a set, we say that $\catC$ is \emph{small}. If $\catC$ is not small we say that it is \emph{large}.
\end{definition}

It is important to note that small categories are necessarily locally small, but the converse does not hold in general.

\begin{proposition}
    Let $\catC$ be a category. The following are equivalent:
    \begin{enumerate}
        \item $\catC$ is small,
        \item $\catC$ is locally small and $\ob \catC$ is a set.
    \end{enumerate}
\end{proposition}
\begin{proof}
    First, suppose $\catC$ is small, i.e. $\mor \catC$ is a set. Then for all $x,y \in \ob \catC$ we have
    \[\Hom_\catC(x,y) \subseteq \mor \catC\]
    and since $\mor \catC$ is a set it follows that $\Hom_\catC(x,y)$ is a set\footnote{exercise!}. Thus, $\catC$ is locally small.

    Next, since $\mor \catC$ is a set, we have that
    \[I \coloneq \{f \in \mor \catC : \exists x \in \ob \catC \,(f = \id x)\} \subseteq \mor \catC\]
    is a set\footnote{by the same exercise!}. Now
    \[\ob \catC = \{\dom f : f \in I\}\]
    and $\{\dom f : f \in I\}$ is a set\footnote{by ZFC's \emph{axiom schema of replacement}}.

    In the other direction, suppose $\catC$ is locally small and that $\ob \catC$ is a set. Now\footnote{by the \emph{union} axiom of ZFC, applied to the family produced by the axiom schema of replacement} we obtain the set
    \[\bigcup_{x \in \ob \catC} \bigcup_{y \in \ob \catC} \Hom_\catC(x,y),\]
    which equals $\mor \catC$.
\end{proof}

This allows us to produce the category $\Set$ of all sets, the category $\Grp$ of all groups, etc. We may also produce the category $\Cat$ of all \emph{small} categories! That's because the data of a small category consists of sets and functions between them, so there is a class of all small categories. On the other hand, we may \emph{not} produce the category of all categories, since the data of an arbitrary large category consists of proper classes and functions between them, and so there is no class of large categories (the elements of a class must be sets!).

\section{Universes}

Classes are somewhat elegant and useful, but also a bit troublesome: sometimes we \emph{would} like to have access to the category of all categories! For example, it would be convenient to have a functor which sends a category $\catC$ to its presheaf category $\Fun(\catC^{\op}, \Set)$. Even if $\catC$ is a small category, its presheaf category can be large (e.g. take $\catC$ to be the terminal category)! The formalism of classes allows us to construct large collections of sets, but it doesn't allow us to construct (even small) collections of classes. To obviate this difficulty, Grothendieck introduced the idea of \emph{universes}, which allow us to filter the class of all sets into layers which we might call ``small sets'', ``large sets'', ``huge sets'', etc. Each of these layers will be a set, i.e. there will be a set of all small sets, and a set of all large sets, and a set of all huge sets, etc!

To be more precise: Grothendieck defined a notion of \emph{universe}, which is just a condition any given set may or may not satisfy. If a set $U$ is a universe, then we should think about $U$ as ``the set of all $U$-small sets''. If we have two universes $U_1$ and $U_2$ with $U_1 \subsetneq U_2$, then we might (for convenience) call the elements of $U_1$ ``small sets'' and the elements of $U_2$ ``large sets''. If we have three universes, ... you get the idea.

Then if we have a universe $U_1$, we can define a ``$U_1$-small category'' to be a category whose class of morphisms is an element of $U_1$ (and in particular is a set). If $U_1 \subsetneq U_2$, it will turn out that the class of all $U_1$-small categories is an element of $U_2 - U_1$, and we will be able to define the category of all $U_1$-small categories, which will itself be $U_2$-small (but not $U_1$-small). And so on, and so forth. Importantly, no matter how many nested universes we have, each universe is a \emph{set}, and so we can deal with its elements and subsets and so on with ordinary set-theory, without ever mentioning classes!

The final philosophical point to make, before we actually define universes, is that at the end of the day, it will always suffice to prove our theorems in some sufficiently large universe -- it is never actually necessary to consider the entire class of sets all at once. After all, any particular mathematical object has a fixed size. So, the plan for setting up category theory with Grothendieck universes is to \emph{only} use small categories (as defined in the previous subsection), but to treat certain sets as ``small'' (these small sets will include e.g. all topological spaces of interest to us in some mathematical context) and other sets as ``large'' (these large sets will include e.g. the set of all small sets).

\begin{definition}
    A set $U$ is said to be a \emph{(Grothendieck) universe} if
    \begin{enumerate}
        \item $U$ is transitive, i.e.
              \[\forall x \forall y (x \in y \in U \implies x \in U),\]
        \item $U$ is closed under forming pairs, i.e.
              \[\forall x,y \in U (\{x,y\} \in U),\]
        \item $U$ is closed under forming power sets, i.e.
              \[\forall x \in U (\mathcal{P}(x) \in U),\]
              where $\mathcal{P}(x)$ denotes the power set of $x$,
        \item $U$ is closed under unions, i.e.
              \[\forall I \in U \,\forall f : I \to U \left(\bigcup_{i \in I} f(i) \in U\right),\]
        \item $\omega \in U$, where $\omega = \{0,1,2,\dots\}$ is the set of finite (von Neumann) ordinals.
    \end{enumerate}
\end{definition}

There are other equivalent ways to express these axioms, but the point is that when $U$ is a universe in ZFC set theory, $(U,\in)$ is a model of ZFC set theory. Some authors (including Grothendieck) do not require the last axiom (in which case $(U,\in)$ will be a model of ZFC minus the axiom of infinity). Another way to say this is that for any universe $U$, we can perform any set-theoretic construction with elements of $U$ and remain inside $U$.

\begin{definition}
    Let $U$ be a universe. By a ``$U$-small set'' we will mean simply an element of $U$. By a ``$U$-small category'' we will mean simply a small category $\catC$ such that $\ob \catC \in U$ and $\mor \catC \in U$.
\end{definition}

We note that there is a set of all $U$-small sets (namely, $U$), and there is a set of all $U$-small categories, but of course neither the set of all $U$-small sets nor the set of all $U$-small categories will be $U$-small.

It is worth pointing out that universes are quite large: every set you've ever seen is $U$-small for \emph{every} universe $U$. That's because every set you've ever explicitly constructed has been built from some iterated application of basic set theory constructions, beginning with the empty set and the set $\omega$, and $U$ is closed under all such constructions\footnote{The only missing step here is to show that $\varnothing \in U$ for all universes $U$: exercise!}! Basic set theory considerations tell us that there are sets much larger than any set you've ever explicitly constructed, and Russell's paradox tells us that we cannot hope to deal with the enormity of the collection of absolutely all sets, but the point of universes is that we don't have to: we can instead restrict our attention to a universe $U$, which will contain every set we care about, and simply work within that universe.

There is only one problem: it is impossible to prove that universes exist\footnote{Unless ZFC is inconsistent, of course. More precisely, ZFC plus ``there exists at least one universe'' proves that ZFC is consistent.}. Instead, we add a new axiom to our set theory which imposes their existence.

\begin{axiom}[Axiom of Universes]\label{ax:universes}
    For all sets $x$, there exists a universe $U$ such that $x \in U$.
\end{axiom}

It is possible in principle that the universe axiom is inconsistent with the other axioms of ZFC, but this is widely believed not to be the case. This belief is not weakly held; set theorists have actually spent a lot of time thinking about this axiom, in the following guise:

\begin{proposition}
    In ZFC, the following two statements are equivalent:
    \begin{enumerate}
        \item The universe axiom holds,
        \item For all cardinals $\kappa$, there exists a strongly inaccessible\footnote{A cardinal $\mu$ is said to be \emph{strongly inaccessible} if it is uncountable, every unbounded subset of $\mu$ has cardinality $\mu$, and $2^\nu < \mu$ for every cardinal $\nu < \mu$.} cardinal $\lambda$ such that $\lambda > \kappa$.
    \end{enumerate}
\end{proposition}

Statement (2) above is an example of what set theorists call a \emph{large cardinal axiom}, which are an important topic in modern set theory. In fact, it turns out that universes are in bijective correspondence with strongly inaccessible cardinals -- for every strongly inaccessible cardinal $\lambda$, the von Neumann stage $V_\lambda$ is a universe, and every universe is equal to $V_\lambda$ for some strongly inaccessible cardinal $\lambda$\footnote{Precisely, for $\lambda$ equal to the cardinality of the universe.}. To summarize: \textbf{the best experts believe strongly that the universe axiom is not inconsistent with the other axioms of ZFC}, but it is impossible to prove this using the axioms of ZFC unless ZFC is itself inconsistent!

Some mathematicians tend to feel uncomfortable with this situation. This fear comes from a good place, but if you find yourself unsettled by this, please note that one could raise the exact same objection about (e.g.) the power set axiom, which simply asserts that every set has a power set. Just as we should be happy to assume (axiomatically) that power sets exist, we should be happy to assume (axiomatically) that universes exist.

In practice, you will often find technical setups of the following form:

\begin{convention}\label{conv:universe}
    Fix universes $U_1$ and $U_2$ with $U_1 \in U_2$\footnote{The universe axiom says in particular that we may make such a choice. It follows from $U_1 \in U_2$ that $U_1 \subsetneq U_2$, and in fact $U_1 \subsetneq U_2$ is equivalent to $U_1 \in U_2$ for any universes $U_1, U_2$.}. Let $\Set$ denote the category of $U_1$-small sets, and let $\SET$ denote the category of $U_2$-small sets. Let $\Cat$ denote the category of $U_1$-small categories, and let $\CAT$ denote the category of $U_2$-small categories.
\end{convention}

Then we have (for example) that $\Set$ is an object of $\CAT$ (exercise!), and we operate \emph{as if} $\Set$ is the category of all sets and $\CAT$ is the category of all categories, even though this is not literally true. Indeed, all of the categories $\Set$, $\SET$, $\Cat$, $\CAT$ described in \Cref{conv:universe} are small in the sense that the class of morphisms in each is a set.

So, our notation here conflicts with our earlier choice to use $\Set$ to denote the category of \emph{all} sets. Besides this kind of notational conflict, the theory of classes does not interfere with the theory of universes in any way -- i.e. it is perfectly fine to discuss proper classes in a setting where one has assumed the universe axiom.

We note now that this setup allows us to define the free cocompletion functor
\[Y : \Cat \to \CAT,\]
which sends a $U_1$-small category $\catC \in \ob \Cat$ to its category of presheaves $\Fun(\catC^\op, \Set)$, and which sends a functor $f : \catC \to \catD$ to left Kan extension along $f^\op$. Note that $Y(\catC)$ is $U_2$-small but typically not $U_1$-small, so $Y$ really does need $\CAT$ (rather than $\Cat$) as its target. We should also be clear that $Y(\catC)$ is the free \emph{$U_1$-cocompletion} of $\catC$, i.e. $Y(\catC)$ admits all $U_1$-small colimits, and all objects of $Y(\catC)$ are $U_1$-small colimits of representable presheaves. Since $\Cat$ is a full subcategory of $\CAT$, we may now compare $Y$ with the inclusion $i : \Cat \hookrightarrow \CAT$, and thereby express the Yoneda embedding as a single natural transformation -- a statement which we could not even write down using classes alone.

\begin{proposition}
    Let $i : \Cat \hookrightarrow \CAT$ denote the inclusion of the full subcategory of $U_1$-small categories. The Yoneda embeddings\footnote{One should of course carefully define the Yoneda embedding and check that it is well-defined: since $\catC$ is $U_1$-small, each hom-set $\Hom_\catC(x,y)$ is indeed an object of $\Set$.}
    \[y_\catC : \catC \to \Fun(\catC^\op, \Set) = Y(\catC)\]
    are the components of a natural transformation $y : i \Rightarrow Y$.
\end{proposition}

The Yoneda lemma itself (the family of bijections $\Hom_{Y(\catC)}(y_\catC(x), F) \cong F(x)$) can be packaged in the same spirit, as a natural isomorphism between two functors defined on a suitable category of triples $(\catC, x, F)$ with $\catC \in \ob \Cat$, $x \in \ob \catC$, and $F \in \ob Y(\catC)$.

\begin{slogan}
    Size issues are unavoidable. However you choose to deal with them, you must always pay attention to which things are large and which things are small.
\end{slogan}
